\section{Autenticación y multi-tenancy}
\label{sec:seguridad}
\textit{Responsable: Mariano (AUTH-01).}

\subsection{Implementación y decisiones}
El stack \texttt{infra/auth} del repositorio \texttt{nubeduca-backend} crea un User Pool de Amazon
Cognito con dos atributos personalizados: \texttt{custom:tenant\_id} (\texttt{utec}, \texttt{upc} o
\texttt{pucp}) y \texttt{custom:role} (\texttt{STUDENT}, \texttt{TEACHER} o \texttt{ADMIN}). Solo un
administrador crea usuarios, y el \textit{app client} solo puede escribir \texttt{email}: un usuario no
puede cambiar su propio tenant ni su rol. Exporta \texttt{nubeduca-<stage>-user-pool-id},
\texttt{-user-pool-arn} y \texttt{-user-pool-client-id}.

Se usa el \textbf{ID token}, porque el access token no trae los atributos personalizados. Cada API
protegida usa un authorizer \texttt{COGNITO\_USER\_POOLS} de API Gateway, que acepta el header
\texttt{Authorization} con o sin el prefijo \texttt{Bearer}. Sin token o con un token inválido,
API Gateway responde 401 antes de invocar la Lambda.

En la plantilla de servicio el authorizer depende del parámetro \texttt{authMode}:
\begin{itemize}
  \item \texttt{authMode=cognito}: las rutas usan el authorizer con el User Pool del stage
        (\texttt{Fn::ImportValue: nubeduca-<stage>-user-pool-arn}). La Lambda lee el tenant y el rol de
        \texttt{requestContext.authorizer.claims}.
  \item \texttt{authMode=mock} (por defecto): sin authorizer y con un token
        \texttt{mock.<tenant>.<ROL>.<userId>}, para el frontend con datos simulados, LocalStack y las
        pruebas.
\end{itemize}
\texttt{/health} queda sin authorizer, porque lo usa la prueba de humo del despliegue.

\paragraph{Aislamiento por tenant.}
Se usa el modelo \textit{pool} (ADR-005): los tres tenants comparten cada tabla, con
\texttt{tenant\_id} como clave de partición. La Lambda toma el tenant \textbf{solo del token}, nunca de
la ruta ni del cuerpo, y todas las consultas van a la partición de ese tenant. El tenant también viaja
en la \texttt{metadata} de cada evento.

\subsection{Resultados y evidencias}
Probado el 21-Nov-2026 en la cuenta de Dev 2 (stage \texttt{dev}) con la plantilla desplegada con
\texttt{authMode=cognito} y los ID tokens reales de los usuarios demo:

\begin{table}[htbp]
\centering\small
\begin{tabularx}{\textwidth}{Xl}
\toprule
Prueba & Resultado \\
\midrule
\texttt{GET /examples} sin token & 401 \\
Token mock con \texttt{authMode=cognito} & 401 \\
Docente UTEC, \texttt{POST /examples} (con \texttt{Bearer}) & 201; \texttt{createdBy} = \texttt{sub} de Cognito \\
Docente UTEC, \texttt{GET /examples} (sin \texttt{Bearer}) & 200, ve su ejemplo \\
Docente UPC, \texttt{GET /examples} & 200 con lista vacía: no ve los datos de UTEC \\
Redeploy con \texttt{authMode=mock} + token mock & 200; método con autorización \texttt{NONE} \\
\bottomrule
\end{tabularx}
\caption{Pruebas del authorizer de Cognito y del aislamiento por tenant.}
\label{tab:auth}
\end{table}

Usuarios demo (creados con \texttt{scripts/crear\_usuarios\_demo.sh}):
\texttt{docente.demo@utec.edu.pe}, \texttt{alumno.demo@utec.edu.pe},
\texttt{docente.demo@upc.edu.pe} y \texttt{alumno.demo@upc.edu.pe}.

\pendiente{Cuando los seis microservicios se desplieguen con \texttt{authMode=cognito} en la cuenta de
integración, agregar una captura del login real en el frontend (sección~\ref{sec:frontend}). No
incluir contraseñas ni tokens en las capturas.}
